\chapter{WASP guess transfer (menu 18)}
\label{ch:menu18}
\menulabel{18}

\section{Purpose}

Starting every structure of a series from an orbital set that already lies in
the right basin -- the other half of the cross-structure problem menu~\menu{17}
maps.  The WASP protocol builds the guess for a target geometry as the
$1/d$-weighted interpolation of the neighbouring structures' orbital
coefficients (the review's Eqs.~(11)--(13)):

\[
\mathbf{C} = \frac{\sum_\beta w_\beta \mathbf{C}_\beta}{\sum_\beta w_\beta},
\qquad w_\beta = \frac{1}{d(\mathbf{G}, \mathbf{G}_\beta)},
\]

over the neighbourhood $d \le \delta$ (RMSD distances), followed by an
orthonormalisation in the target geometry's overlap metric.  This menu writes
the result into a \emph{gbw-ready} Molekel mkl: \code{orca\_2mkl <name>.fbk
-gbw} converts it into a \code{.gbw} that ORCA reads through \code{!moread} +
\code{\%moinp}.

\section{What you need}

The neighbour exports (localized orbitals of the already-computed structures),
the target's own export (its overlap metric orthonormalises the mixture) and
the target's mkl template (\code{orca\_2mkl <target> -mkl} after any cheap run
at that geometry).  All structures share the basis set, the atom order and the
orientation; the manifest schema is in Chapter~\ref{ch:formats}.

\section{Walkthrough}

\lstinputlisting[style=fbkcode, firstline=54, lastline=72,
  title={The menu-18 example session (the frozen N$_2$ scan: the two far
  neighbours build the guess for the middle geometry)}]
  {examples/expected/m18_guess.txt}

\section{What you get}

The weights table, the orthonormalisation residual, the \code{.fbk.mkl} file
(the exact conversion and \code{\%moinp} command are printed with it) and a
report (\file{<manifest>.guess.fbk.md}) with the citations.  The occupations
and orbital energies written along come from the nearest neighbour.

\section{How to read the numbers}

\begin{readbox}
\begin{itemize}
  \item \textbf{The weights are $1/d$ over the neighbourhood} -- closer
    structures dominate, and a neighbour \emph{at} the target geometry takes
    everything (the weight would diverge): the guess then is that structure's
    own orbital set, and the report says so.
  \item \textbf{The residual is the certificate of the write}: the mixture is
    orthonormal in the target's overlap metric before it is written
    ($3\times10^{-15}$ in the example).  The write-back route itself was
    measured on ORCA 6.1.1: the converted file is accepted as
    \code{INITIAL GUESS: MOREAD}, and a gbw$\to$mkl$\to$gbw round trip
    preserves the coefficients to the file's print precision ($5\times10^{-8}$).
  \item \textbf{A guess selects the basin, not just the starting point.}  In
    the frozen example the interpolated guess converged the target geometry's
    RHF to a \emph{second} solution 19~mEh \emph{below} the branch the scan
    itself used -- a measured multiple-solution instance (RHF/def2-SVP at
    1.6~\AA).  After any guess transfer, check that the solution still belongs
    to the series (energy, orbital character, the mapping); a good guess does
    not guarantee branch continuity.
\end{itemize}
\end{readbox}

\section{Boundaries and related sections}

\begin{refusalbox}
\begin{itemize}
  \item Structures with different atoms, basis or orbital counts are refused;
    a delta that excludes every neighbour is refused with the closest
    distance named.
  \item Nothing is aligned silently: the RMSD (and the coefficient
    interpolation behind it) assumes the structures share one frame, which is
    how a scan is normally produced.
  \item The written file is a guess: ORCA re-optimises the orbitals, and this
    menu claims nothing about convergence.
\end{itemize}
\end{refusalbox}

Related: \menu{17} (the mapping that keeps the active space consistent -- the
guess and the map are complements), \menu{12} (the entropy selection the map
carries), and Chapter~\ref{ch:criteria}.
